\begin{table}[t]
\centering
\caption{Standardized Effect Sizes}
\label{tab:sde}
\begin{adjustbox}{max width=\textwidth}
\begin{tabular}{lcccccc}
\toprule
Outcome & $\hat{\beta}$ & SE & SD($Y$) & SDE & SE(SDE) & Classification \\
\midrule
\addlinespace
\multicolumn{7}{l}{\textit{Panel A: Pooled}} \\
\addlinespace
Log firm births & -0.2152 & 0.0908 & 1.250 & -0.1722 & 0.0727 & Large negative \\
Birth rate (\%) & -1.9728 & 0.8396 & 4.039 & -0.4884 & 0.2078 & Large negative \\
Log active enterprises & -0.0263 & 0.0377 & 1.221 & -0.0215 & 0.0309 & Small negative \\
\addlinespace
\multicolumn{7}{l}{\textit{Panel B: Heterogeneous (sample splits)}} \\
\addlinespace
Log births (regulated sectors) & -0.2342 & 0.0995 & 1.014 & -0.2310 & 0.0982 & Large negative \\
Log births (non-regulated sectors) & -0.1773 & 0.0810 & 1.150 & -0.1542 & 0.0705 & Large negative \\
\bottomrule
\end{tabular}
\end{adjustbox}
\par\vspace{0.3em}{\footnotesize
\textit{Notes:} \textbf{Country:} European Union (27 member states). \textbf{Research question:} Does late transposition of EU directives into national law suppress firm entry in affected sectors? \textbf{Policy mechanism:} EU directives set a transposition deadline by which member states must enact national legislation; when the deadline passes without transposition, firms face regulatory uncertainty about which rules apply, raising the option value of delaying entry. \textbf{Outcome definition:} Log of firm births (new enterprise registrations) from Eurostat Business Demography, measured at the country--NACE-section--year level; birth rate is births per 100 active enterprises. \textbf{Treatment:} Binary indicator equal to one when at least one directive affecting a sector has passed its transposition deadline without the member state notifying national measures. \textbf{Data:} CELLAR SPARQL (directive transposition records) and Eurostat bd\_enace2\_r3 (business demography), 2008--2022, country $\times$ NACE section $\times$ year. \textbf{Method:} Two-way fixed effects with country $\times$ sector and year fixed effects; standard errors clustered at the country level; wild cluster bootstrap (Webb weights, 9,999 iterations). \textbf{Sample:} EU-27 member states, NACE sections mapped to directive subject areas via keyword classification of EUR-Lex titles. SDE $= \hat{\beta} / \text{SD}(Y)$ where SD($Y$) is the pre-treatment standard deviation. Classification refers to magnitude, not statistical significance: Large ($|$SDE$| > 0.15$), Moderate ($0.05$--$0.15$), Small ($0.005$--$0.05$), Null ($< 0.005$).
}
\end{table}
